% Capitolo 3 - Riforme strutturali per l'adeguatezza nel sistema NDC
% Ultima modifica: 2026-05-28
% Fonti principali usate: MEF_RGS2025, MEF_RGS2025_NdA, Mazzaferro2023,
%   FrancoTommasino2020, Gronchi2024, BevilacquaGronchi2025, Altiparmakov2022,
%   Marchionne2004, OCPI2026, RosenGayerRapallini2023, BosiGuerra2022,
%   Bosi2018, Artoni2014

%% ───────────────────────────────────────────────────────────────────
\chapter{Riforme strutturali per l'adeguatezza nel sistema NDC}
\label{ch:riforme}
%% ───────────────────────────────────────────────────────────────────

% --- §3.1 La domanda che la riforma deve risolvere ---------------------
\noindent
La diagnosi del capitolo precedente lascia in eredità un paradosso operativo.
Il sistema pensionistico italiano è contabilmente sostenibile in tutti gli
scenari macroeconomici prevedibili dalla Ragioneria Generale dello Stato e
dall'Ageing Working Group della Commissione europea \cite{MEF_RGS2025,
MEF_RGS2025_NdA, FrancoTommasino2020}.\footnote{Per scenario \emph{prevedibile}
si intende qui l'insieme delle proiezioni elaborate da RGS e dall'Ageing Working
Group con le rispettive analisi di sensitività su fecondità, longevità,
produttività e flussi migratori, non qualunque scenario immaginabile. La
sostenibilità contabile è garantita dentro questo perimetro tecnico: la critica
dell'osservatorio sui conti pubblici discussa più avanti agisce sui margini di
quelle ipotesi, non al loro esterno.} Il paradosso è operativo, non retorico:
la traiettoria della spesa in rapporto al PIL regge anche nelle sensitività
simmetriche su quei parametri, ma la sostenibilità è acquistata comprimendo il
tasso di sostituzione delle generazioni che si pensionano interamente sotto il
regime contributivo. Il sistema regge perché scarica sui lavoratori giovani
l'aggiustamento che le coorti più anziane non hanno sopportato.

\medskip
\noindent
La domanda di questo capitolo è di conseguenza una sola: \emph{come si
riporta il tasso di sostituzione verso una soglia di adeguatezza senza
rompere la sostenibilità contabile, e a quale costo, ripartito su quali
generazioni}? Le riforme che la letteratura italiana e internazionale propone
si dividono in quattro famiglie: parametriche dentro l'NDC, redistributive
sulla perequazione, strutturali sull'architettura a pilastri, contabili sulla
separazione fra previdenza e assistenza. Ognuna ha un costo in punti di PIL,
una platea di beneficiari e una fattibilità politica diversa. La tesi che il
capitolo sostiene è che le riforme percorribili sono quelle che riducono il
trasferimento implicito di rischio dalle coorti anziane a quelle giovani senza
ricorrere a una transizione di sistema, perché la transizione di sistema, dove
è stata tentata in carve-out, è fallita \cite{Altiparmakov2022}.

% --- §3.2 Le leve di policy e cosa ciascuna risolve --------------------
\section{Le leve di policy e cosa ciascuna risolve}
\label{sec:leve}

\noindent
Il manuale italiano di scienza delle finanze \cite{Bosi2018, BosiGuerra2022,
RosenGayerRapallini2023} riconduce gli interventi sul sistema previdenziale
a sei leve, derivabili dalla decomposizione canonica del rapporto fra spesa
pensionistica e PIL: età legale di pensionamento, immigrazione, tasso di
fecondità, importo medio della pensione, aliquota contributiva, ricorso al
debito pubblico. Non tutte agiscono sulle stesse variabili. Una distinzione
operativa fra leve che migliorano la sostenibilità e leve che migliorano
l'adeguatezza chiarisce dove sia possibile lavorare (tabella~\ref{tab:leve}).

\begin{table}[h!]
\centering
\small
\renewcommand{\arraystretch}{1.35}
\caption{Le sei leve di policy e gli effetti attesi su sostenibilità e
         adeguatezza}
\label{tab:leve}
\begin{tabularx}{\textwidth}{X p{0.22\textwidth} p{0.22\textwidth} p{0.16\textwidth}}
\toprule
\textbf{Leva} & \textbf{Sostenibilità (Sp/PIL)} & \textbf{Adeguatezza (TS)}
              & \textbf{Costo politico} \\
\midrule
Età legale di pensionamento $\uparrow$ & migliora & migliora via montante & alto \\
\midrule
Immigrazione netta $\uparrow$ & migliora & neutrale & alto \\
\midrule
Fecondità $\uparrow$ & migliora (lungo termine) & neutrale & nessuna leva diretta \\
\midrule
Importo pensione $\downarrow$ & migliora & peggiora & alto \\
\midrule
Aliquota contributiva $\uparrow$ & migliora & migliora (NDC) & medio \\
\midrule
Ricorso al debito pubblico $\uparrow$ & neutrale nel breve; rinvia il costo come debito (anche implicito) & migliora (transitorio) & basso \\
\bottomrule
\end{tabularx}
\par\smallskip
\footnotesize Fonte: rielaborazione da \cite{Bosi2018} e
\cite{RosenGayerRapallini2023}; la lettura del ricorso al debito come spesa
previdenziale implicita si rifà a \cite{Fornero2018}. La colonna costo politico
riflette su quale coorte grava l'onere, secondo il criterio dell'elettore
mediano esposto in chiusura della sezione.
\end{table}

\noindent
Le sei leve non agiscono sullo stesso punto del rapporto. Età legale e importo
medio incidono sul numeratore, la prima riducendo gli anni di erogazione e
accrescendo il montante, il secondo tagliando la prestazione; fecondità,
immigrazione e aliquota agiscono invece sulla base che lo alimenta, il monte
contributivo e il PIL a cui è agganciato. Che il riferimento sia il PIL e non i
soli contributi non è dettaglio contabile: nulla impone che la previdenza sia
finanziata solo dai contributi sociali, e destinarvi gettito di altra natura è
scelta politica, non vincolo giuridico. Le tre leve sulla base si distinguono però
per orizzonte: la fecondità conta solo nel lunghissimo periodo, perché un nuovo
nato versa contributi dopo vent'anni; l'immigrazione netta migliora il rapporto
subito ma rinvia il problema, perché chi versa oggi maturerà domani un proprio
diritto, e se i flussi si interrompono è la coorte autoctona a finanziare le
prestazioni di chi può aver lasciato il paese. È lì che il patto
intergenerazionale, sostenibile sulla carta, diventa politicamente fragile.

\medskip
\noindent
Il ricorso al debito appartiene a una categoria diversa dalle prime cinque.
Finanziare a debito una quota di spesa non riduce il rapporto spesa/PIL nei conti,
e nel breve non lo peggiora: ne sposta il costo nel tempo. Dietro c'è il debito
pensionistico implicito, il valore attuale delle prestazioni promesse al netto dei
contributi futuri attesi, che un sistema a ripartizione accumula per costruzione e
che la tradizione italiana di Castellino e Fornero legge come la componente
nascosta della spesa previdenziale \cite{FrancoMarinoZotteri2004,
CastellinoFornero1997, Fornero2018}. Quanto pesi è controverso, perché dipende dal
tasso di sconto adottato e un livello anche elevato resta compatibile con una
ripartizione sostenibile \cite{FrancoMarinoZotteri2004, BeltramettiDellaValle2011}.
Il punto qualitativo regge comunque: il debito non riequilibra il sistema, lo
rinvia, a costo politico basso oggi e a carico delle coorti future.

\medskip
\noindent
Il problema italiano non è mancanza di leve ma eccesso di vincoli politici. Le tre leve a costo
elettorale alto, età, immigrazione, riduzione importo, sono quelle che già
operano automaticamente attraverso i meccanismi di aggiustamento del metodo
contributivo \cite{MEF_RGS2025_NdA}. L'aliquota contributiva è già al 33\%
per i dipendenti, fra le più elevate dell'area OCSE \cite{OCPI2026,
BosiGuerra2022}, e ulteriori aumenti aggraverebbero il cuneo fiscale già in
cima alla classifica europea. Il debito pubblico, fra i più elevati dell'area euro in rapporto al PIL, non consente espansioni straordinarie. Le riforme percorribili sono
dunque quelle che lavorano \emph{dentro} il sistema NDC, ridistribuendo il
rischio fra coorti anziché aumentando la spesa aggregata.

\medskip
\noindent
Che un vincolo politico sia stringente o cedevole non è giudizio
impressionistico ma ha una struttura analitica, ed è la stessa che gradua le
colonne di costo politico e fattibilità di questo capitolo. La teoria
dell'economia politica della previdenza la deriva dal teorema dell'elettore
mediano: in una democrazia la dimensione di un sistema a ripartizione è fissata
dalla coorte decisiva, e \citet{Browning1975} mostra che quella coorte, in
prevalenza di mezza età e anziana, sostiene un sistema più grande di quello
attuarialmente efficiente, perché incassa i benefici delle prestazioni che la
riguardano e scarica sui contribuenti futuri la quota di costo che non la tocca.
La rassegna di \citet{GalassoProfeta2002} generalizza il punto: una riforma che
riduce le promesse pensionistiche è tanto più difficile da approvare quanto più
la coorte mediana è vicina alla quiescenza, perché le si chiede un sacrificio i
cui benefici maturano per chi voterà fra trent'anni. \citet{SinnUebelmesser2003}
ne offrono la versione operativa adottata qui: distinguono l'\emph{età di
indifferenza}, la coorte che da una riforma non guadagna né perde, dall'età
della coorte politicamente decisiva, e osservano che quando la seconda supera la
prima la maggioranza capace di approvare la correzione cessa di esistere.
L'Italia, con un'età mediana dell'elettorato fra le più alte d'Europa e una
platea di percettori di pensione in crescita, si colloca dal lato sfavorevole di
quella soglia. È questo criterio, e non la cronaca del dibattito, a graduare la
fattibilità nelle pagine che seguono: una riforma è politicamente praticabile
quando il suo costo ricade su coorti future o non ancora decisive, oppure è
differito nel tempo; è impraticabile quando intacca la posizione della coorte
che oggi decide; sta nel mezzo quando richiede un trasferimento esplicito che
compete, nel bilancio corrente, con la spesa che quella stessa coorte presidia.

% --- §3.3 Riforme parametriche dentro l'NDC -----------------------------
\section{Riforme parametriche dentro l'NDC}
\label{sec:parametriche}

\subsection{Ricalibrazione dei coefficienti di trasformazione}
\label{sub:coefficienti}

\noindent
Il coefficiente di trasformazione è il meccanismo attraverso cui il sistema
NDC assorbe l'aumento della speranza di vita. Per il biennio 2025-2026 il
coefficiente a 67 anni è fissato al 5,608\%, e le proiezioni della RGS
indicano la sua progressiva discesa man mano che la speranza di vita residua
aumenta di circa 4,8 anni per gli uomini e 4,1 per le donne entro il 2070
\cite{MEF_RGS2025}. La revisione biennale ha funzionato come ammortizzatore
automatico della longevità e concorre alla riduzione cumulata di oltre
sessanta punti percentuali di PIL garantita dal ciclo di riforme avviato nel
2004; di quella riduzione la Ragioneria attribuisce più di un terzo
all'adeguamento automatico dei requisiti anagrafici di accesso \cite{MEF_RGS2025_NdA}.

\medskip
\noindent
Il dispositivo presenta però una criticità tecnica documentata. Sandro
Gronchi, coautore della riforma Dini, mostra che le tavole di sopravvivenza
utilizzate per costruire i coefficienti italiani fanno riferimento a coorti
che sono in media quindici anni e mezzo più vecchie di quelle che entrano in
pensione, producendo una sopravvalutazione sistematica del coefficiente e, di
conseguenza, \guillemotleft{}quote di pensione superiori ai contributi versati che, da un
lato, si traducono in doni ai pensionati, mentre dall'altro producono sbilanci
fra la spesa pensionistica e il gettito contributivo \guillemotright{}\cite{Gronchi2024}.
La distanza tra la coorte di riferimento (\emph{mutuante}), le cui tavole di
sopravvivenza alimentano il coefficiente, e la coorte che va effettivamente in
pensione e a cui il coefficiente viene applicato (\emph{mutuataria}), sale a
diciannove anni quando il pensionamento avviene a sessanta anni. Il dispositivo dei
coefficienti, pur robusto in aggregato, premia sistematicamente le pensioni
anticipate.

\medskip
\noindent
La proposta di Gronchi consiste nel sostituire le pensioni anticipate con un
\emph{assegno di accompagnamento} sul modello svedese, erogato fra i
sessantatré e i sessantasei anni come prestito a titolo oneroso, rimborsato
al raggiungimento dell'età di vecchiaia con un prelievo sulla pensione
definitiva \cite{Gronchi2024}. Il meccanismo conserva la flessibilità ed
elimina il trasferimento implicito a carico della fiscalità generale. Gronchi
non ne stima il costo aggregato, e una stima ufficiale non esiste; due dati
verificati ne fissano però l'ordine di grandezza. Nel 2022 le pensioni anticipate
hanno sfiorato il 62\% delle liquidazioni ai dipendenti \cite{Gronchi2024}, e le
sole deroghe \emph{Quota 100}, \emph{102} e \emph{103} fra il 2019 e il 2024 sono
costate oltre nove miliardi, circa un miliardo e mezzo l'anno, per una platea molto
più ristretta \cite{MEF_RGS2025_NdA}. Se una finestra così circoscritta pesa già a
quel ritmo, il trasferimento implicito sull'intera platea delle anticipate è
verosimilmente di alcuni miliardi annui. La cifra resta indicativa; la sproporzione
che la genera no, ed è il punto: il dispositivo dei coefficienti sussidia con la
fiscalità generale la maggioranza delle nuove liquidazioni ai dipendenti.

\subsection{Pensione contributiva di garanzia per le carriere discontinue}
\label{sub:pcg}

\noindent
L'NDC nella sua formulazione pura non garantisce un minimo pensionistico, e
questo non è un difetto del modello ma un suo principio di disegno
\cite{FrancoTommasino2020}. Daniele Franco e Pietro Tommasino lo riconoscono
esplicitamente quando osservano che \guillemotleft{}an NDC scheme cannot guarantee a
minimum pension income \guillemotright{}\cite{FrancoTommasino2020}, e proprio per questo
sostengono che la previdenza pubblica non debba essere assistenza, due
funzioni che andrebbero finanziate da fonti diverse. Il problema operativo è
che la generazione che entra oggi nel mercato del lavoro presenta carriere
strutturalmente discontinue, e per una quota non trascurabile di lavoratori
atipici, lavoratrici con percorsi interrotti da maternità o caregiving,
autonomi a basso reddito, il montante NDC alla fine della vita attiva sarà
insufficiente persino a raggiungere la soglia minima di accesso al
pensionamento anticipato, fissata dal 2030 in tre volte l'assegno sociale:
poiché nel 2025 l'assegno sociale vale circa 539 euro mensili, poco più di
7.000 euro l'anno \cite{INPSAssegnoSociale2025}, la soglia corrisponde a una
pensione di circa 21.000 euro annui \cite{MEF_RGS2025_NdA}.

\medskip
\noindent
La proposta di una \emph{pensione contributiva di garanzia}, formulata e
valutata per il caso italiano da \citet{Raitano2020}, colma la
distanza fra la pensione NDC calcolata sui contributi effettivi e una soglia di
adeguatezza ancorata a un multiplo dell'assegno sociale, attraverso un
trasferimento esplicito finanziato dalla fiscalità generale. Il meccanismo si
distingue dall'attuale integrazione al trattamento minimo per tre ragioni.
La fonte di finanziamento è dichiarata, e quindi sottoposta a vincolo di
bilancio trasparente. La platea è definita ex ante, perché l'integrazione
scatta solo per chi resta sotto la soglia di garanzia. La logica contributiva
del sistema NDC è preservata, perché chi ha versato abbastanza non riceve
nulla di aggiuntivo. Il costo annuo a regime, sulla base delle proiezioni RGS
sulla composizione delle nuove pensioni \cite{MEF_RGS2025}, è stimabile fra
tre e cinque miliardi: una cifra significativa ma inferiore al costo cumulato
delle deroghe \emph{Quota 100}, \emph{102} e \emph{103} fra il 2019 e il 2024,
che secondo le valutazioni RGS ha superato i nove miliardi
\cite{MEF_RGS2025_NdA, BosiGuerra2022}.

\subsection{Il paradosso italiano dell'indicizzazione}
\label{sub:indicizzazione}

\noindent
Mirko Bevilacqua e Sandro Gronchi documentano una distorsione tecnica trascurata
dal dibattito \cite{BevilacquaGronchi2025}. La riforma Dini capitalizza i montanti
al PIL nominale, ma rivaluta le pensioni già in pagamento solo ai prezzi al
consumo, non al PIL al netto di uno sconto come la regola NDC richiederebbe. Il
coefficiente di trasformazione fissa però l'assegno iniziale \emph{come se} questo,
una volta liquidato, dovesse crescere ogni anno un punto e mezzo meno del PIL: uno
sconto che nei fatti non viene mai applicato. Il risultato è una pensione iniziale
più alta di quanto il sistema potrebbe sostenere, il cui potere d'acquisto si erode
poi lungo tutta la fase di pagamento. Gli autori lo chiamano \emph{paradosso
italiano}, e ne attribuiscono la sopravvivenza al \guillemotleft{}duplice scudo
dell'insostenibilità sociale delle vie d'uscita e dell'indisponibilità della classe
politica a impegnarsi nella faticosa comprensione dei tecnicismi
\guillemotright{}\cite{BevilacquaGronchi2025}.

\medskip
\noindent
Le vie d'uscita che propongono si distinguono per lo sconto applicato
all'indicizzazione al PIL, e nessuna è oggi sul tavolo. Agganciare le pensioni in
essere al PIL meno 1,5 punti preserva l'assegno iniziale ma ne erode il potere
d'acquisto, severamente nel periodo 2010-2026; lo sconto norvegese di 0,75 punti
riduce l'erosione abbassando un po' la pensione iniziale; l'indicizzazione piena
protegge il potere d'acquisto ma taglia le pensioni iniziali del cinque o sei per
cento, con un effetto politico esplosivo. La sola correzione che il dispositivo
accetterà senza traumi è quella già incorporata nella revisione biennale dei
coefficienti.

% --- §3.4 Riforma della perequazione --------------------------------
\section{La riforma della perequazione legittimata dalla Corte}
\label{sec:perequazione}

\noindent
La perequazione delle pensioni in essere, la rivalutazione annuale degli assegni
in pagamento per adeguarli al costo della vita, è la voce più rigida del bilancio
INPS, ma anche quella su cui dopo la sentenza della Corte costituzionale 167/2025,
depositata il tredici novembre 2025, si è aperto lo spazio politico più ampio
\cite{CorteCost167_2025}. La Corte ha statuito che \guillemotleft{}non sussiste un
imperativo costituzionale che imponga l'adeguamento annuale di tutti i trattamenti
pensionistici \guillemotright{}, legittimando la pratica già adottata dal governo nel
2023 di differenziare le aliquote di rivalutazione per scaglione, dal 32 all'85\%
sopra quattro volte il minimo. Una perequazione strutturalmente selettiva sulle
pensioni alte può liberare risorse dell'ordine di diversi miliardi l'anno a regime,
in misura che dipende dalle soglie e dalle aliquote scelte: è un potenziale che le
simulazioni circolate nel dibattito pubblico collocano fra i più alti di qualunque
intervento sulla spesa pensionistica \cite{EconomiaItalia2026_video}.

\medskip
\noindent
La leva ha due caratteristiche che la distinguono dalle altre. Non richiede
trasferimenti aggiuntivi a carico della fiscalità generale, perché libera risorse
dentro il perimetro INPS; e non tocca le pensioni dei più deboli, perché scatta
solo sopra una soglia multipla dell'assegno sociale. La sua difficoltà è politica,
non tecnica: il principio dell'invarianza nominale delle pensioni è protetto dalla
coalizione elettorale più ampia di qualunque altra politica di spesa. È la sentenza
della Corte a cambiare le condizioni di fattibilità, perché avendo escluso un
obbligo di adeguamento integrale rende difendibile una perequazione differenziata
prima esposta al rischio di incostituzionalità. Resta una misura redistributiva ex
post, estranea alla logica di corrispettività che governa il resto del capitolo: la
si accoglie non perché attuarialmente neutra, ma perché libera risorse per i
trasferimenti che la neutralità attuariale, da sola, non finanzia.

% --- §3.5 Carve-out o add-on: la lezione est-europea ---------------------
\section{Architettura a pilastri: carve-out o add-on?}
\label{sec:carveout}

\noindent
La proposta storica di riforma strutturale è quella che \citet{Marchionne2004}
avanzò su \emph{Moneta e Credito}: un sistema misto a tre pilastri, con il pubblico
a ripartizione ad aliquota ridotta al 25\%, un secondo pilastro occupazionale a
capitalizzazione all'8,5\% e un terzo individuale al 4\%. La difesa della
sostenibilità poggiava sull'idea che dirottare contributi verso la
capitalizzazione, mantenendo lo stesso \emph{replacement ratio} con aliquote più
basse, liberi le risorse per finanziare la transizione \cite{Marchionne2004}. È una
privatizzazione a \emph{carve-out}, dove parte dei contributi PAYG va a fondi a
capitalizzazione obbligatoria.

\medskip
\noindent
L'argomento ha un punto debole preciso. La transizione impone alla generazione di
passaggio un doppio onere: pagare le pensioni in essere, che nessuno ha
capitalizzato, e insieme accantonare il proprio risparmio. Marchionne lo vede e
ammette che il vecchio onere non può essere coperto per intero in quel modo, ma lo
tratta come gestibile, affidandosi alla condizione, mai quantificata, che i
rendimenti della capitalizzazione superino in valore attuale i costi del
transitorio. Le sue simulazioni misurano i tassi di sostituzione a regime, non il
fabbisogno fiscale della transizione, ed è lì che il costo viene sottostimato.

\medskip
\noindent
La proposta non è mai stata implementata in Italia. Lo è stata, fra il 1998
e il 2008, in dieci paesi dell'Europa orientale e, nel 2022, in Grecia. Il
working paper di Nikola Altiparmakov per il CeRP raccoglie l'evidenza
empirica del primo quarto di secolo dell'esperimento \cite{Altiparmakov2022}.
Il rendimento reale medio dei fondi privati obbligatori nei dieci paesi dell'est è
stato del 2,2\% annuo fra il 1998 e il 2020, contro una crescita del PIL del 2,7\%,
con un differenziale negativo di 0,6 punti rispetto al rendimento implicito del PAYG
sostituito \cite{Altiparmakov2022}. I costi di transizione, proiettati nell'ordine
di pochi miliardi, si sono rivelati di un ordine di grandezza superiore: l'autorità
attuariale greca stima la transizione in cinquanta-settanta miliardi su cinquant'anni
contro i sei del piano iniziale. Altiparmakov conclude che il \emph{carve-out}
\guillemotleft{}has been thoroughly debunked \guillemotright{}dai dati e ne identifica il
meccanismo finanziario come \guillemotleft{}the root cause of reform reversals in Eastern
Europe \guillemotright{}\cite{Altiparmakov2022}. Ungheria, Polonia, Estonia e Slovacchia
sono tornate indietro fra il 2010 e il 2020 sotto la pressione di crisi fiscali.

\medskip
\noindent
La lezione operativa che la tesi accoglie è netta. Una transizione di
sistema finanziata sottraendo contributi al primo pilastro non è un'opzione
percorribile per l'Italia, che ha un debito pubblico maggiore di quello dei
paesi dell'Europa orientale al momento delle loro privatizzazioni e una
crescita potenziale strutturalmente inferiore. Il secondo pilastro va
costruito come \emph{add-on}, attraverso il conferimento volontario del TFR
e la deducibilità fiscale fino a 5.164,57 euro annui dei contributi
aggiuntivi \cite{MEF_RGS2025}. Il capitolo successivo, nella parte dedicata alla
previdenza complementare, approfondisce le ragioni per cui l'add-on, oggi sotto-utilizzato
(le risorse del secondo pilastro italiano valgono il 9,5\% del PIL, una
frazione dei multipli del prodotto raggiunti nei paesi a forte capitalizzazione
come Paesi Bassi e Danimarca, \citealp{FrancoTommasino2020}), è una
riforma fattibile a costo politico basso e a impatto attuariale neutrale.

% --- §3.6 Separazione previdenza/assistenza -----------------------------
\section{Separare previdenza e assistenza}
\label{sec:separazione}

\noindent
La leva contabile più trascurata dal dibattito previdenziale italiano è la
separazione fra le voci di spesa propriamente previdenziali, finanziate dai
contributi, e le voci assistenziali finanziate dalla fiscalità generale.
L'osservatorio sui conti pubblici stima che i trasferimenti dello Stato
all'INPS per oneri pensionistici siano stati 97 miliardi nel 2024, di cui
36 miliardi imputabili a oneri previdenziali specifici \cite{OCPI2026}. La
distinzione è operativa: la prima cifra include integrazioni al minimo,
maggiorazioni sociali, pensioni di invalidità civile, mentre la seconda
riflette interventi su specifiche gestioni in disavanzo. La spesa per
prestazioni assistenziali coperte da fiscalità generale ammontava già nel
1997 al 4\% del totale delle pensioni \cite{Milazzo1999}, e la commistione
ha origini istituzionali precise nella disciplina degli anni Settanta.

\medskip
\noindent
La separazione contabile produrrebbe tre effetti misurabili. Il primo è
maggiore trasparenza sul tasso di rendimento implicito del sistema
contributivo, perché permetterebbe di calcolare quanto del 33\% di aliquota
finanzi prestazioni assicurative e quanto finanzi sussidi assistenziali. Il
secondo è una possibile riduzione dell'aliquota contributiva dell'ordine di
tre o quattro punti percentuali, se le componenti assistenziali fossero
trasferite alla fiscalità generale. Il terzo è di natura politica: rende
esplicito chi paga e a beneficio di chi, sottraendo all'INPS la funzione di
camera di compensazione fra responsabilità previdenziali e politiche sociali.
Il vincolo è di natura aggregata, perché la separazione richiederebbe il
finanziamento esplicito, attraverso imposte dirette o indirette, di una spesa
che il Centro Studi e Ricerche Itinerari Previdenziali colloca nell'ordine dei
quaranta miliardi annui: nel 2023 le componenti assistenziali contabilizzate
dentro la spesa pensionistica, dalle maggiorazioni sociali e dalle integrazioni
al minimo fino alle prestazioni interamente assistenziali come invalidità civile,
indennità di accompagnamento e assegni sociali, sfioravano i quarantasei miliardi
\cite{ItinerariPrevidenziali2024}. Senza questa precondizione la separazione
resta esercizio contabile.

% --- §3.7 Il costo di sospendere gli automatismi -----------------------
\section{Il costo di sospendere gli adeguamenti automatici}
\label{sec:nonriformare}

\noindent
L'argomento più forte a favore del mantenimento degli adeguamenti automatici
già in vigore, i coefficienti di trasformazione e i requisiti anagrafici di
accesso, è quantificato dallo stesso MEF nello scenario controfattuale della
loro sospensione. La nota
di aggiornamento al Rapporto n.~26 stima che il blocco simultaneo dei due
meccanismi a partire dal 2025 comporterebbe un incremento del rapporto fra
debito pubblico e PIL di cinquantotto punti percentuali al 2070
\cite{MEF_RGS2025_NdA}: più di mezzo punto di PIL all'anno, in media, di
debito accumulato per non lasciare che i meccanismi tecnici facciano il loro
lavoro. La sola sospensione dei requisiti anagrafici costerebbe trentasei
punti, quella dei soli coefficienti ventidue.

\medskip
\noindent
La cifra ha due implicazioni. Smontare le riforme parametriche già in essere
è quantitativamente più costoso di qualunque riforma proattiva qui discussa.
Le leve che operano oggi automaticamente, longevity link sui coefficienti,
aggancio biennale dell'età alla speranza di vita, sono la condizione minima
di permanenza della sostenibilità, non un punto di partenza per
ammorbidimenti. Carlo Mazzaferro quantifica il costo storico del difetto
opposto, l'eccessiva generosità della transizione 1995-2011, in ottanta
miliardi di euro di passività previdenziali implicite, pari al 5,5\% del PIL
italiano \cite{Mazzaferro2023}. La cifra misura il sussidio implicito
trasferito alle coorti che hanno mantenuto il regime retributivo nonostante
l'introduzione del metodo contributivo, e ricade per costruzione sulle coorti
successive. Il 92,6\% degli individui pensionati fra il 1996 e il 2019, nel
campione amministrativo INPS analizzato da Mazzaferro, presenta un \emph{Net
Present Value Ratio} maggiore di uno, ovvero riceve in valore attuale più di
quanto ha versato \cite{Mazzaferro2023}. Non si tratta di una previsione: si
tratta del passato già metabolizzato dal sistema.

\medskip
\noindent
Il dato di Mazzaferro va incrociato con la critica dell'osservatorio sui conti
pubblici alle proiezioni RGS \cite{OCPI2026}. Il rientro della spesa dopo il picco
dipende dall'ipotesi che la produttività cresca all'1,04\% medio annuo per
quarantacinque anni, contro lo 0,12\% registrato fra il 2015 e il 2025. Non è un
numero arbitrario, perché discende dalle ipotesi armonizzate dell'Ageing Working
Group europeo \cite{MEF_RGS2025}; la fragilità sta nella distanza dal dato recente,
depresso da doppia recessione e pandemia, mentre anche la media italiana di lungo
periodo resta sotto quella soglia. Se il rientro non si materializza, la
sostenibilità formale regge solo perché il rischio è già scaricato sui pensionati
futuri attraverso i coefficienti. La sostenibilità italiana non è un equilibrio ma
una tregua condizionata, e lo stesso MEF lo riconosce quando scrive che
\guillemotleft{}la sostenibilità finanziaria, perseguita senza tenere conto di un
adeguato livello delle prestazioni, potrebbe determinare complessi effetti sociali,
rendendo inevitabili successive revisioni del quadro normativo \guillemotright{}
\cite{MEF_RGS2025}.

% --- §3.8 Una matrice operativa delle riforme -------------------------
\section{Quali riforme sono percorribili}
\label{sec:matrice}

\noindent
La tabella \ref{tab:matrice} sintetizza le riforme discusse lungo quattro
dimensioni: l'efficacia rispetto all'adeguatezza, misurata in punti di tasso di
sostituzione recuperabili; il costo in punti di PIL su orizzonte di lungo
periodo; la fattibilità politica; e l'allineamento con il principio di
neutralità attuariale che la tesi assume come bussola \cite{FrancoTommasino2020,
Gronchi2024}.

\begin{table}[h!]
\centering
\small
\renewcommand{\arraystretch}{1.35}
\caption{Riforme strutturali del primo pilastro: una matrice operativa}
\label{tab:matrice}
\begin{tabularx}{\textwidth}{X X p{0.14\textwidth} p{0.12\textwidth} p{0.14\textwidth}}
\toprule
\textbf{Riforma} & \textbf{Effetto TS} & \textbf{Costo (PIL)}
                 & \textbf{Fattibilità} & \textbf{Neutralità att.} \\
\midrule
Assegno di accompagnamento (Gronchi)
   & marginale & $-0,2\%$ & medio-alta & sì \\
\midrule
PCG per carriere discontinue
   & da 0 a $+30$ pp\,$^{a}$ & $+0,15$ a $+0,25\%$ & media & sì \\
\midrule
Indicizzazione coerente al PIL
   & $-3$ pp iniziali, $+x$ in vita & $\approx 0$ & bassa & sì \\
\midrule
Perequazione selettiva pensioni alte
   & nullo & $\approx -0,5\%$\,$^{b}$ & medio-alta & no \\
\midrule
Carve-out misto (Marchionne)
   & incerto & $+0,8$ a $+1,7\%$\,$^{c}$ & bassa & sì \\
\midrule
Add-on volontario (TFR + ded. fiscale)
   & da $+5$ a $+10$ pp & $-0,3\%$ (mancato gettito) & alta & sì \\
\midrule
Separazione previdenza/assistenza
   & indiretto & neutrale a regime & bassa & sì \\
\bottomrule
\end{tabularx}
\par\smallskip
\footnotesize
Fattibilità graduata secondo il criterio dell'elettore mediano discusso nella
sezione~\ref{sec:leve} \citep{Browning1975, GalassoProfeta2002,
SinnUebelmesser2003}: alta quando il costo ricade su coorti future o non
decisive, bassa quando colpisce la coorte vested che oggi decide.\\
$^{a}$ Solo per la platea sotto la soglia di garanzia, stimabile in oltre un
milione di unità a regime \cite{MEF_RGS2025_NdA}.\\
$^{b}$ Risparmio indicativo, dipendente da soglie e aliquote; fattibilità
sostenuta dalla legittimazione della Corte costituzionale 167/2025.\\
$^{c}$ Range stimato dall'evidenza est-europea su orizzonte cumulato di
cinquant'anni \cite{Altiparmakov2022}.
\end{table}

\noindent
La matrice mostra una distribuzione delle riforme che la tesi non può
chiudere in un ranking univoco. Le riforme con la fattibilità politica più
alta, add-on e perequazione selettiva, hanno effetti collaterali opposti
sull'adeguatezza: l'add-on la migliora ma solo per chi vi aderisce
attivamente, e l'adesione attuale è ferma a circa un lavoratore su tre
\cite{FrancoTommasino2020}; la perequazione selettiva non tocca l'adeguatezza dei
nuovi pensionati ma libera risorse spendibili per i più deboli. Le riforme
con la neutralità attuariale più pura, assegno di accompagnamento e PCG,
hanno fattibilità media perché richiedono un trasferimento esplicito dalla
fiscalità generale alle gestioni INPS. Le riforme strutturali che la dottrina
italiana ha sostenuto per due decenni, sintetizzate dalla proposta Marchionne,
sono quelle che l'evidenza empirica est-europea ha falsificato, e la tesi le
considera non percorribili nella loro forma carve-out. La via che resta è una
combinazione: rafforzamento del secondo pilastro volontario come add-on per
chi è in grado di aderire, pensione contributiva di garanzia come pavimento
per chi non lo è, perequazione selettiva come fonte di finanziamento per i
trasferimenti necessari.

\medskip
\noindent
Una leva ulteriore esula dalla matrice, perché non ritocca un parametro ma erode il
denominatore. Il monte contributivo è reddito da lavoro: se l'automazione sposta
valore aggiunto dal lavoro al capitale, la base si assottiglia a prescindere dai
coefficienti, e gli stessi parametri rendono meno gettito. Non è un'ipotesi di
scuola, perché la diffusione dei robot comprime occupazione e salari nei mercati
esposti \cite{AcemogluRestrepo2020} e l'automazione cognitiva estende l'effetto alle
professioni a media qualificazione \cite{AcemogluJohnson2023}. Il contributo
sull'automazione proposto in Senato appartiene a questa famiglia, con un gettito
stimato attorno agli 8 miliardi annui per un fondo a sostegno dei lavoratori
sostituiti e della previdenza \cite{Bacchiocchi2025}. La sua efficacia è però
incerta: la letteratura sull'imposizione ottimale trova che l'imposta sui robot
efficiente è piccola e dà guadagni modesti, perché gran parte del beneficio passa
dalla riforma del sistema fiscale nel suo complesso \cite{Thuemmel2022}, e tassare il
capitale automatizzato per finanziare il welfare rischia di abbassare occupazione e
crescita senza rimpiazzare il gettito perso dal lavoro \cite{Furgase2025}. Il punto
tocca un nodo di questo capitolo: l'automazione è proprio la leva che potrebbe
colmare il divario di produttività su cui poggia il rientro della spesa, e tassarla
rischia di indebolire la condizione che dovrebbe difendere. Se il gettito non
sparisce ma si sposta verso il capitale, la risposta coerente è allargare la base
imponibile in quella direzione dentro una struttura fiscale riformata, non un
tributo isolato sui robot \cite{Hotte2024}.

\medskip
\noindent
Il capitolo successivo passa dalla riforma di sistema alla sua versione
individuale. La domanda diventa quella che un lavoratore della Generazione Z
si pone già oggi: data l'età legale di pensionamento crescente, dato il
tasso di sostituzione strutturalmente in discesa, dato il rischio politico
che le riforme qui discusse non vengano implementate, quali sono le
strategie operative che un singolo lavoratore può mettere in campo per
costruire da solo l'adeguatezza che il sistema pubblico non garantirà più?
Il problema si sposta dal disegno istituzionale alla pianificazione
finanziaria personale, ma la materia sottostante è la stessa: un sistema
previdenziale che funziona come tregua e non come equilibrio, e che richiede
a ogni coorte che entra di costruirsi un margine di sicurezza.
